\subsection{Veronese像の全排除と零交換子CYの必要幾何被覆}
\label{subsec:cubic-CY-Veronese}
\begin{proposition}[固定as4入力の必要部分]\label{prop:cubic-CY-Veronese-excluded}
CY AS4型14641の二次関係空間を $R$ とし、
$R\cap\bigwedge^2V=0$ とする。
その対称射影は三次元部分空間 $W\subset V$ に対する $\operatorname{Sym}^2W$ にならない。
\end{proposition}
\begin{proof}
固定出典\cite{AS4CYpinned}の
\texttt{CY\_ZERO\_COMMUTATORS\_VERONESE\_EXCLUSION.tex}の排除定理を、必要性だけ用いる。
ここではその全被覆と当repoで再検証した入力を明記する。

$W=\langle a,b,c\rangle$ とする。
六Veronese二次式の全21積のrankは15で、二次関係は対称 $3\times3$ 行列の余因子で表せる。
係数行列のrank1,2,3は関係のrank3,4,6に対応するので、rank6の対称topは
一つの可逆行列 $A$（$\det A=-1/4$）へ正規化できる。
全巡回条件の混合核は25次元、残る交代--交代部分は $C_0+\gamma J$ であり、
source rank6の条件と実関係は
\[
 C_0+\gamma J-B^tA^{-1}B=0,\qquad R=U+KB^tA^{-1}
\]
となる。$U,K$ は元の対称・交代tensor基底である。
$d$ を含む三列 $B_d$ は、全三次式 $P(a,b,c)$ の偏微分 $\tfrac13\partial P$ に一致する。
その像は $A^{-1}$ に関して等方であり、最大等方三平面に含まれる。

固定出典の同変等長写像
$\bigwedge^2(\operatorname{Sym}^3k^2)\simeq\operatorname{Sym}^2W$
の全係数を検証した。
最大等方平面の二種類 $s\wedge k^4$ と $\bigwedge^2(s^\perp)$、および二変数三次式 $s$ の
三重根・二重根・三単根の三型から、全六平面を得る。
各平面に入る全偏微分という線型条件を十個の三次係数のまま解き、全核を確認する。
$u=a+ib$, $v=c$, $w=-a+ib$ とすれば、
全非零三次式は平行移動・尺度変更により
\[
 u^3,\quad u^2v,\quad u^2(-w),\quad u^2(-w+u),\quad
 u(-uw+8v^2),\quad u(-uw+8v^2)+u^3,\quad P_\alpha,\quad P_\beta
\]
の八位置で覆われる。
最後の二つは三単根の二種類の一次元核であり、全式を入力JSONに保持する。
$P=0$ は別に扱い、$\gamma=0$、全実関係が $W^{\otimes2}$ に入るため、
自由積によって三次元成分が少なくとも25となり排除される。

八位置の全線型解と全残余idealを元係数から再構成すると、最初の四位置は実平方
$(a+ib)^2$ を関係に含み、次の二位置は単位idealとなる。
最後の二位置はそれぞれ両符号の $\gamma=i/3$、$\gamma=3\sqrt3/4$ と三自由係数である。
$d$ の全 $W$ 方向への平行移動はこの三係数全体へ移り、Jacobian行列式は
$8/27$、$-72$ で非零である。全16語の実関係が移ったことも直接確認した。
対角変換 $\operatorname{diag}(1,-1,-1,-1)$ は両符号を交換する。
従って有限の二代表だけで全候補を扱える。
元六関係から当repoの独立な商計算を行うと、それぞれ
\[
 (\dim C_0,\dim C_1,\dim C_2,\dim C_3)=(1,4,10,22),\qquad
 (\dim C_0,\ldots,\dim C_4)=(1,4,10,20,37)
\]
となり、AS4 Hilbert値に反する。
この排除は全最大等方平面・全三次境界・全自由係数・両符号を含む。
\end{proof}

\begin{theorem}[零交換子CYの必要完全交差像]\label{thm:cubic-CY-zero-commutator-CI22}
H3のCY二次包絡で $R_2\cap\bigwedge^2V=0$ なら、その六次元対称射影の像は
整な三次元完全交差 $CI(2,2)\subset\mathbb P^5$ である。
これは必要な対称射影像の被覆であり、その上の全包絡・切断の分類やAS十分性を主張しない。
\end{theorem}
\begin{proof}
二次関係空間を $R_2$、対称射影を $U$ とする。
$c=0$ なので $\dim U=6$ であり、$R_2$ は $U$ 上のgraphである。
対称なmiddle top行列の像が $R_2$ であることから、$SS$ 部分もrank6である。
全対称化が零なので、それは六二次式間のrank6関係 $Q$ を与える。

像の次元を $r\le2$ と仮定する。
ある線型切断 $\mathbb P^r\subset\mathbb P^3$ への制限は像上優越である。
実際、微分の一般rankは $r$ であり、ある定義域内の点でその補空間を接空間に持つ線型切断を選べる。
像は非退化なので、この制限の二次式空間は六次元のままである。
$r\le1$ は $h^0(\mathcal O_{\mathbb P^r}(2))<6$ に反する。
$r=2$ なら平面全二次式系となり、像はVeronese面である。

その標的座標で、元の六二次式を成分とする対称 $3\times3$ 行列はrank1となる。
分数体上で $Z=\rho ff^t$ と書き、$f$ を共通因子のない多項式vectorにする。
各既約因子の付値を $Z_{ii}=\rho f_i^2$ で比較すれば $\rho$ は多項式である。
六式の共通因子は無い（共通一次因子を持つ二次式空間は高々四次元）ので $\rho$ は単元となる。
斉次次数から $f_i$ は一次式であり、六式の独立性から三つは独立である。
従って元の $U$ 自体が $\operatorname{Sym}^2W$ となり、前命題に反する。
よって像の次元は3である。

命題\ref{prop:cubic-CY-six-quadrics-second}から独立な第二二次関係を得る。
滑らかな $Q$ の錐の環は、座標を選んで
\[
 k[z_0,\ldots,z_5]/(z_0z_5-z_1^2-z_2^2-z_3^2-z_4^2)
\]
と書け、一意分解整域である。
具体的にはこれは正常環で、$z_0$ は素元、その可逆化はLaurent多項式環である。
素因子をこの可逆化から降ろすと、追加される高さ1素因子は主因子 $(z_0)$ だけなので、因子類群は零となる。
像の錐の素idealはこの環内で高さ1ゆえ、一つの斉次式で生成される。
第二二次関係が存在するのでその次数は高々2、非退化性から1ではない。
従って像はスキームとして整な $CI(2,2)$ となる。
\end{proof}

独立再現器は\\
\texttt{verify\_cubic\_CY\_Veronese.py}。
固定上流証明書から最小入力を保持し、当repoで25混合核の完全性、六平面の全三次核、
全八ideal、全平行移動と両符号、二代表の元商を独立計算する。
上流全検証器の再実行と、この独立再生は区別して記録する。
幾何学的被覆の通常証明は固定出典および上の証明に依存し、形式検証や新規性確認とはしない。
次には $CI(2,2)$ の七射影型と全表示への必要被覆を監査する必要がある。
